\setcounter{aufgabe}{28}
\hypertarget{e4}{}\einheitenkopf{Einheit 4 von 5 · Nullstellen und Schnittpunkte der Parabel}

\begin{aufgabe}{Nullstellen -- Was wird gleichgesetzt? Schreib nur die Gleichung hin,
die du aufstellen musst. Löse nichts.}
\begin{teile}
\teil Nullstellen von $p(x) = x^2 - 5x + 4$ \hfill \leerfeld
\teil Schnittpunkte von $p(x) = x^2 + 2x$ und $f(x) = 3x + 6$ \hfill \leerfeld
\teil Stellen, an denen $p(x) = x^2 - 4x$ den Wert $12$ hat \hfill \leerfeld
\teil Nullstellen von $p(x) = (x - 7)^2 - 9$ \hfill \leerfeld
\teil Schnittpunkte von $p(x) = (x + 1)^2$ und $f(x) = -x + 5$ \hfill \leerfeld
\end{teile}
\end{aufgabe}

\begin{aufgabe}{Nullstellen -- aus der Scheitelpunktform. Setze $p(x) = 0$ und zieh die
Wurzel. Gib die Nullstellen an.}
\beispiel{(x - 4)^2 - 9 &= 0 &&\mid + 9 \\
(x - 4)^2 &= 9 \\
x - 4 &= 3 && \text{oder}\quad x - 4 = -3 \\
x_1 &= 7 && x_2 = 1}
\begin{gleichungsraster}[4]
\gl{p(x) = (x - 3)^2 - 16} & \gl{p(x) = (x + 2)^2 - 25} \\
\gl{p(x) = (x - 5)^2 - 4} & \gl{p(x) = (x + 6)^2 - 1} \\
\gl{p(x) = (x - 8)^2} & \gl{p(x) = (x - 2)^2 + 5} \\
\gl{p(x) = -(x - 1)^2 + 9} & \\
\end{gleichungsraster}
Sieh nur auf den Scheitel und die Öffnung. Wie viele Nullstellen hat die Parabel?
Begründe.
\begin{teile}
\setcounter{teil}{7}
\teil $p(x) = (x + 4)^2 + 2$
\teil $p(x) = -(x - 3)^2 + 7$
\end{teile}
\end{aufgabe}

\begin{aufgabe}{Nullstellen -- aus der Normalform. Setze $p(x) = 0$ und rechne mit der
Lösungsformel.}
\beispiel{x^2 - 2x - 15 &= 0 \\
x_{1,2} &= 1 \pm \sqrt{1 + 15} = 1 \pm 4 \\
x_1 &= 5 && x_2 = -3}
\begin{gleichungsraster}[4]
\gl{p(x) = x^2 - 7x + 10} & \gl{p(x) = x^2 + 9x + 20} \\
\gl{p(x) = x^2 - 3x - 18} & \gl{p(x) = x^2 + 2x - 24} \\
\end{gleichungsraster}
Teile zuerst durch die Vorzahl von $x^2$.
\begin{gleichungsraster}[5]
\gl{p(x) = 2x^2 - 10x + 12} & \gl{p(x) = 3x^2 + 3x - 18} \\
\end{gleichungsraster}
Geht die Wurzel nicht auf, gib die Nullstellen zuerst mit Wurzel und dann als
Näherungswert auf zwei Stellen an.
\begin{gleichungsraster}[5]
\gl{p(x) = x^2 - 4x + 1} & \gl{p(x) = x^2 - 8x + 14} \\
\end{gleichungsraster}
\hfill (P10 2025)
\end{aufgabe}

\begin{aufgabe}{Nullstellen -- am Graphen ablesen. Gezeichnet sind zwei Graphen: der
gekrümmte gehört zur Parabel $p$, der gerade zur Geraden $f$.}

\begin{ksys}[xmin=-3,xmax=6,ymin=-5,ymax=6,ablesen]
\parabel{1}{1}{-4}{}
\gerade{3}{-7}{}
\end{ksys}

\begin{teile}
\teil Lies die beiden Nullstellen von $p$ ab. \hfill \pfeld{N_1} \ \pfeld{N_2}
\teil Lies die beiden Schnittpunkte von $p$ und $f$ ab. \hfill \pfeld{S_1} \ \pfeld{S_2}
\teil Jan sagt: „Jede Gerade schneidet die Parabel in zwei Punkten.`` Beurteile die
Aussage. \hfill (P10 2023)
\end{teile}
\end{aufgabe}

\begin{aufgabe}{Parabel -- $x$-Wert zum $y$-Wert. Berechne, für welche $x$ die Parabel
den angegebenen Wert annimmt.}
\begin{teile}
\teil $p(x) = x^2 + 4x$, gesucht sind die Stellen mit $p(x) = 12$.
\teil $p(x) = x^2 - 6x$, gesucht sind die Stellen mit $p(x) = 7$.
\teil $p(x) = -x^2 - 2x + 7$, gesucht sind die Stellen mit $p(x) = -8$.
\hfill (P10 2023)
\end{teile}
\end{aufgabe}

\begin{aufgabe}{Schnittpunkte -- rechnen. Setze die Terme gleich, löse die Gleichung
und berechne die $y$-Werte mit der Geraden. Gib die Schnittpunkte als Punkte an.}
\beispiel{x^2 + 3x + 1 &= x + 9 \\
x^2 + 2x - 8 &= 0 \\
x_1 &= 2 && x_2 = -4 \\
S_1(2 \mid 11) && S_2(-4 \mid 5)}
\begin{teile}
\teil $p(x) = x^2 + 5x$ \quad und \quad $f(x) = 2x + 4$
\teil $p(x) = x^2 - 4$ \quad und \quad $f(x) = 3x + 6$
\teil $p(x) = x^2 + 2x + 3$ \quad und \quad $f(x) = 6x + 3$
\end{teile}
Die Parabel steht in der Scheitelpunktform. Löse zuerst die Klammer auf.
\begin{teile}
\setcounter{teil}{3}
\teil $p(x) = (x - 2)^2 + 2$ \quad und \quad $f(x) = 2x + 1$
\teil $p(x) = (x - 4)^2 - 5$ \quad und \quad $f(x) = -3x + 11$ \hfill (P10 2022)
\end{teile}
Hier sind es zwei Parabeln.
\begin{teile}
\setcounter{teil}{5}
\teil $p(x) = x^2 + x - 3$ \quad und \quad $q(x) = 2x^2 - 5x + 5$
\end{teile}
\end{aufgabe}

\begin{aufgabe}{Schnittpunkte -- Punktprobe. Setze den Punkt in beide Funktionen ein
und antworte.}
\begin{teile}
\teil Zeige, dass $P(3 \mid 4)$ auf $p(x) = x^2 - 5$ und auf $f(x) = 2x - 2$ liegt.
\teil Liegt $Q(-2 \mid 1)$ auf $p(x) = x^2 - 3$ und auf $f(x) = -x - 1$?
\teil Liegt $R(1 \mid 0)$ auf $p(x) = x^2 + 3x - 4$ und auf $f(x) = 2x - 1$?
\hfill (P10 2014)
\end{teile}
\end{aufgabe}

\begin{aufgabe}{Schnittpunkte -- Fehler finden. Mia sucht die Schnittpunkte von
$p(x) = x^2 + 4x$ und $f(x) = x + 4$ und rechnet:}
\rechnung{x^2 + 4x &= x + 4 \\
x^2 + 3x - 4 &= 0 \\
x_1 &= 1 && x_2 = -4}
Sie schreibt als Antwort $S_1(1 \mid 0)$ und $S_2(-4 \mid 0)$. Finde den Fehler,
benenne ihn und korrigiere die Antwort.

\begin{teile}
\teil Berechne die Schnittpunkte von $p(x) = x^2 - 2x$ und $f(x) = x + 4$.
\end{teile}
\end{aufgabe}

\begin{aufgabe}{Nullstellen -- Begründen. Antworte ohne lange Rechnung.}
\begin{teile}
\teil Begründe, warum die Nullstellen einer Parabel genau ihre Schnittpunkte mit der
$x$-Achse sind.
\teil $p(x) = x^2 - 6x + 5$ hat den Scheitel $S(3 \mid -4)$. Gib eine waagerechte
Gerade an, die die Parabel nicht schneidet, und begründe.
\teil Begründe, warum eine nach unten geöffnete Parabel, deren Scheitel über der
$x$-Achse liegt, immer zwei Nullstellen hat.
\end{teile}
\end{aufgabe}
